\section{GPU-resident effects: grass and fire}
\label{sec:effects}

\subsection{The shared span}

Vegetation and fire both need information from the chunk's surface, but
neither needs a CPU copy of the triangle mesh.  They consume the same
published span:
\begin{itemize}
\item vegetation reads the vertex and index pools directly;
\item fire reads the node SDFs on the host build and emits anchors.
\end{itemize}
Both are keyed by chunk id and both participate in the chunk's lifecycle:
a republish replaces the old data, a delete tombstones it, and a rebuild
creates a new version.

\subsection{Vegetation placement}

The placement shader \code{VegetationPlacement.comp} reads the chunk's
vertex span, treats every triangle as a potential grass emitter, and:
\begin{enumerate}
\item rejects triangles with no grass-brush vertex;
\item computes the triangle area;
\item expands a stochastic number of virtual slots proportional to area;
\item samples a barycentric point;
\item applies the dominant-vertex brush test;
\item applies the biome noise gate;
\item selects a blade variant and a rotation fraction by position hash;
\item writes a 32-byte instance record.
\end{enumerate}
The placement output goes into a scratch ring.  The host reads the accepted
count, allocates an exact persistent buffer, copies the scratch slice, and
publishes the chunk.  No instance byte reaches the host.  A coarse rung
with thousands of tufts is capped at \code{kPlacementScratchPerChunk} to
bound the far-field draw; the cap is a deliberate density reduction, not a
silent overflow.

\subsection{Why not one instance per triangle?}

The old CPU code expanded slots on the CPU, one per area-weighted sample.
The GPU port does the same expansion in the shader: each invocation is one
triangle and loops over its stochastic slot count.  This avoids a
prefix-sum pass and a second dispatch.  The trade-off is that the output
order is non-deterministic; the instance record itself is deterministic
because every random draw is seeded by the chunk, triangle, and slot.

\subsection{The vegetation draw}

The cull emits a blade-tuft command with
\code{indexCount=36} (12 blades, one triangle each) and an impostor command
with \code{indexCount=6}.  The instance count and first instance come from
the per-chunk table.  The vegetation vertex shader reads the instance
record:
\begin{itemize}
\item position,
\item normal,
\item variant and rotation,
\item baked height,
\end{itemize}
and expands the procedural blade geometry.  The impostor path samples the
captured blade atlas and reconstructs depth from the captured depth array.

\subsection{Lava/fire anchors}

The SDF fire effect is driven by brush-4 surfaces.  The old collector
sampled the CPU mesh; after the GPU-only change it received an empty
geometry.  The replacement resolves anchors during the chunk build:
\begin{enumerate}
\item for every brush-4 surface node that owns a vertex slot, evaluate the
  node's trilinear SDF and gradient at its centre;
\item take one Newton step to the zero crossing;
\item clamp to the node cube;
\item emit a position, normal, and seed.
\end{enumerate}
The anchors are carried on \code{LoDMesh} to the main thread and ingested
by \code{SdfRenderer}.  The renderer caps the number per chunk, deduplicates
nested LoD rungs on an 8 m grid, builds a flame instance list, and stages
the SDF scene.  The volumetric fire path then renders the flames through
the same generic SDF scene as smoke and bullets.

\begin{figure}[htbp]
\centering
\begin{tikzpicture}[font=\small]
  \node[io] (V) {vertex/index span};
  \node[gpu,below=of V] (veg) {placement compute};
  \node[gpu,right=of veg] (sdf) {node SDF $\to$ anchor};
  \node[io,below=of veg] (scratch) {scratch ring};
  \node[io,right=of scratch] (anchors) {anchor list};
  \node[cpu,below=of scratch] (copy) {exact buffer copy\\+ publish};
  \node[io,below=of copy] (draw) {grass draws};
  \node[gpu,below=of anchors] (fire) {volumetric fire};
  \draw[arrow] (V)--(veg); \draw[arrow] (veg)--(scratch);
  \draw[arrow] (scratch)--(copy); \draw[arrow] (copy)--(draw);
  \draw[arrow] (V)--(sdf); \draw[arrow] (sdf)--(anchors);
  \draw[arrow] (anchors)--(fire);
\end{tikzpicture}
\caption{Two GPU-resident effects sharing the chunk span.  Vegetation
writes instances and copies them to an exact persistent buffer; fire
resolves anchors during the build and hands them to the SDF renderer.}
\label{fig:effects}
\end{figure}

\subsection{Ownership and cleanup}

Both effects register with the chunk lifecycle:
\begin{itemize}
\item a republish replaces the chunk's instance buffer/anchor entry;
\item an empty chunk tombstones the old entry;
\item a chunk delete removes the entry and defers the buffer destroy until
  all pending submissions complete;
\item app shutdown stops the worker pools before the device idle wait.
\end{itemize}
The effect data is therefore never left pointing at a freed span.
